\section{Norms of Poisson differences}
\label{p14:comp:poisson-differences}

The following calculation proves Proposition~\ref{p1:discussion:charlier-norm}.
It fixes both the signed-difference convention and the factor of two between
$\ell^1$ distance and total variation.

\begin{proof}[Proof of Proposition~\ref{p1:discussion:charlier-norm}]
Write $v_k=(\Pi_\mu*\Delta^{*2})\{k\}$.  Direct calculation gives
\[
 v_k=\pi_\mu(k)-2\pi_\mu(k-1)+\pi_\mu(k-2)
     =\frac{\pi_\mu(k)}{\mu^2}\bigl((k-\mu)^2-k\bigr).
\]
The quadratic has roots $u_-,u_+$ and is strictly negative
exactly at the integers $a,\ldots,b$.  These form a nonempty
interval; integer roots contribute zero and are excluded by
the definitions of $a,b$.  Since $\sum_kv_k=0$,
\[
 \begin{aligned}
 \sum_k|v_k|
   &=-2\sum_{k=a}^b v_k\\
   &=2\bigl\{\pi_\mu(a-1)-\pi_\mu(a-2)
                 +\pi_\mu(b-1)-\pi_\mu(b)\bigr\},
 \end{aligned}
\]
by telescoping.  This proves \eqref{p1:discussion:finite-norm}.
For $\mu=2$, the quadratic is $(k-1)(k-4)$: the negative
masses are $v_2=-e^{-2}$ and $v_3=-2e^{-2}/3$.
Their sum, multiplied by $-2$, gives
\eqref{p1:discussion:lambda-two}.  The same calculation at
$\mu=1$ gives the value $3/e$ used in
\eqref{p1:close:dyadic-example}.

For the general bound, define the normalized Charlier polynomials
\[
 C_r(k;\mu)=\sum_{j=0}^r
       \binom rj(-1)^{r-j}\frac{(k)_j}{\mu^j}.
\]
Then $(\Pi_\mu*\Delta^{*r})\{k\}=\pi_\mu(k)C_r(k;\mu)$.
Their exponential generating function is
\[
 \sum_{r\ge0}C_r(k;\mu)\frac{t^r}{r!}
       =e^{-t}(1+t/\mu)^k.
\]
For $P\sim\Pois(\mu)$, the expectation of the product of
the generating functions with variables $s,t$ is
\[
 \E\!\left[e^{-s-t}(1+s/\mu)^P(1+t/\mu)^P\right]
       =e^{st/\mu}.
\]
Coefficient comparison, justified in a neighborhood of
$(s,t)=(0,0)$ by the Poisson exponential moments, yields
\[
 \E[C_r(P;\mu)C_\ell(P;\mu)]
       =\ind_{\{r=\ell\}}\,\frac{r!}{\mu^r}.
\]
Cauchy--Schwarz now gives the second bound in
\eqref{p1:discussion:general-norm}.  The first follows from
$\|\Pi_\mu\|_1=1$, $\|\Delta\|_1=2$, and submultiplicativity
of the convolution norm.
\end{proof}
